\chapter[Sermon 27. Here Is Lively Described the Lamb in the Throne of God, Receiving the Book of the Hand of Him That Sits and Opening It.]{Here Is Lively Described the Lamb in the Throne of God, Receiving the Book of the Hand of Him That Sits and Opening It.}
\label{sermon:027}
\markright{Sermon 27. Here Is Lively Described the Lamb in the Throne of ...}

And one of the Elders said to me, “Do not weep.” Behold, the Lion from the tribe of Judah, the root of David, has prevailed to open the book and to break the seven seals thereof. And I looked, and behold, in the midst of the throne, and of the four living creatures, and in the midst of the Elders, stood a lamb as though he had been slain, which had seven horns and seven eyes, which are the seven spirits of God sent into all the world, and he came and took the book out of the right hand of him who sat on the throne.

Since John had wept, that no one was worthy even to look upon the book belonging to the one who sat on the throne, one of the twenty-four elders comforted him. His name is not expressed, and it seems vain and curious to seek it. Nevertheless, some commentators suppose him to be the patriarch Jacob, since shortly after his oracle or prophecy is recited. And so the author proceeds in a most orderly fashion to the description also of the Son of God, by whom the heavenly Father, as all Scripture affirms everywhere, governs all things. Hitherto he has described the one who sits on the throne, and before that, the Holy Spirit. Therefore, these are wholesome and most profitable doctrines for the church, by which the true faith is confirmed.

The comfort of this elder, and truly the heavenly and most godly doctrine aims at this: that we should understand that all the complaints, weeping, grudging, and the various distresses of our mind cannot be quenched, appeased, and quieted unless we see and believe that to Christ (as is most plainly and manifestly set forth here) has been given by the Father all power in heaven and on earth. Therefore, to be constituted both the only Redeemer and also the head, Prince, and governor of all, which under the seal of faith and truth should govern all things that are by God’s providence ordained, and even now despises them, and reveals to us so much of God’s judgments as suffice us. If we believe this with a faithful and sincere mind, we shall have quiet consciences in all the works of God, even those that are hard to endure and seem to some men most unreasonable. For we know that he by whom all things are governed is of our nature and kind, yes even our own brother. And he who truly favors us with all his heart has suffered death for us, and loves nothing better in all the world than humankind. Moreover, he who has overcome death, the Devil, and Hell, and has overcome them for us. Who will now suspect his government, permission, or operation? You have a brother in the Prince’s Court, whom you are assured will favor you from the bottom of his heart. You hear how he has given to him of the Prince the government and judgment of the whole country; would you hesitate or be reluctant to submit yourself to him? No, rather you trust and hope to obtain anything from your brother.

\index{Hebrews, Epistle to the}Therefore let us remember how the Scripture does not teach this only, but everywhere, that Jesus Christ, the Son of God, and in truth of the same substance with us after his humanity, in dying for us, deserved to have a name given him, which is above all names, and that all things should be subject to his rule, whatever may be in the visible or invisible world. For so John testifies in the first chapter. And Paul also to the Philippians, Colossians 1, and to the Hebrews in the first chapter. He is said at this present to have overcome and obtained to open the book and loose the seals thereof. Therefore, by knowledge of him and through faith in him, we obtain that with a joyful mind we may look upon the book, the judgments, and all the works of God, and quietly and patiently bear the opening thereof and the rule of all together.

But so that we may judge more rightly of Christ as ruler of all—although he has already depicted him vividly—now he proceeds to portray him in his most godly and excellent qualities, so that we should not be afraid of his rule, nor unwilling to submit ourselves wholly and quietly to his governance.

First, it is said that a Lion of the tribe of Judah has overcome, namely, that same Christ of old: to have overcome the Devil, sin, death, the world, hell, and all power of the adversary. And he overcame in dying, and so achieved the highest dignity, and was made Lord of all. The Devil is also called a lion by Saint Peter. Solomon and the Prophets call tyrants lions. Our author therefore calls Christ a lion, not of the common sort, but of the tribe of Judah. For he alludes to the prophecy of the patriarch Jacob, which is in Genesis 49. He prophesies there that Shiloh shall come, with plenty and good fortune, which like a lion that has taken its prey, neither is there any man that can drive him from it, can defend those that are his, whom he has caught out of the dragon’s claws, so that no hostile power dares once hiss against him. Christ therefore is declared a victor or conqueror greatest, most mighty, and most invincible. Which belongs to him alone. Yet you shall find kings who are every hour overcome by wicked lusts, which will suffer themselves to be called invincible. Briefly, this first note in the description of Christ shows that Jesus Christ, governor of all, is that very same whom the patriarchs and prophets have prophesied to come into the world, a prince most invincible.

\index{Jerome, Saint}\index{Mahomet}\index{Justification}Secondly, Christ is called the root of David, wherein he appears to have alluded to that saying of Isaiah in the 11th chapter. Then shall a bud come forth of the stock of Jesse, and a flower shall ascend out of the roots thereof. Namely, Mary, the daughter of David, from whom that most sacred flower, Christ, sprang and came, was the stock of Jesse. And of the very roots of David, or of the virgin—I mean of the most true human nature—Jesus Christ was born very man into the world. For he took not the angelic nature, but the seed of Abraham. He is therefore our brother, of the same substance with us, after his humanity. These things do comfort us exceedingly and confute heretics strongly, who deny that Christ has a very human body. We have more hereof in Matthew 1 and Luke 1:2-3. After it is expressly spoken of the same our Lord, that he is in the midst of the throne, in the midst of the four beasts, and in the midst of the twenty-four Elders: and is therefore excepted out of the number of creatures, out of the number of angels, and out of the number of saints. For he is greater than these, to wit, of the same substance with the Father, in glory and power equal. For the Father is in the midst of the throne, from thence proceeds the Holy Ghost: even there is found also now the Lamb, Christ, not only very man, but also very God. And is a distinct person. For the blessed Trinity knows no confusion. The Father is God, the Son is God, the Holy Ghost is God: yet are all three but one God, the Father in his subsistence, the Son in his, and the Holy Ghost in his, not making three Gods, but three properties and persons in one indivisible and eternal essence. And whereas Christ is mentioned to be in the midst of the beasts and in the midst of the Elders, he is doubtless signified after the divine nature to be everywhere, to be the life and preservation of all creatures, also in the midst of his chosen and of his Church. Therefore, like as we believe Jesus Christ to be very man, so let us also believe him to be very God, of the same substance with God the Father. Therefore let Servetus perish with Arius and Mahomet, and as many as deny Christ to be the Son of God, coequal with the Father in all things. Furthermore, he is now also called a Lamb, not that he is a sheep by nature, but for that by a lamb is prefigured the innocent redeemer of the world and the only wholesome sacrifice of all the faithful. A lamb is a token of innocence, and from the beginning appointed for sacrifices. Abel offered up a lamb; after the law was offered a daily sacrifice, in the morning a lamb, and at evening a lamb. For Christ is the expiation of those who were in the beginning of the world and who in the end shall be. The Paschal lamb in Exodus 12 whose blood prohibited the destroying angel from the houses and tents, represented the figure of Christ, by whose precious blood we are reconciled to God. This exposition of the Paschal lamb, Jerome, Peter in 1 Peter 1, and Paul in 1 Corinthians 5, have brought. Isaiah accords with them in the 53rd chapter. And so expounded by the Apostle Philip in Acts 8. Finally, John the Baptist, which with the finger stretched out and pointing to Christ, exclaimed: “Behold the Lamb of God, which takes away the sins of the world.” Let us therefore believe that the same Jesus Christ, unto whom all power is given of the Father, is our deliverer, our expiation, reconciliation, innocence, sanctification, justification and everlasting salvation: as he whom we shall hear in the 13th chapter to have been slain from the beginning of the world, for so much as his only death and one oblation made from the beginning of the world, and continually to the world’s end, doth sanctify all those that are sanctified, which the Apostle also affirms in Hebrews 10.

Nevertheless, this Lamb or Savior of the world is said to stand in the midst of the throne, truly because now he executes the office of a universal king, priest, and governor, always ready and prepared to save. Thus, Saint Stephen also in Acts 7 sees him standing. Or else, in other places, we read that Christ sits on the right hand of the Father. But this passage does not contradict that, considering that to sit signifies both rest and reign.

\index{Augustine, Saint}\index{Resurrection}Moreover, this our Lamb appears on the throne of divine majesty as if he were slain, not because he was not truly slain and dead—for that is described in detail shortly thereafter—but because he did not remain in death, but rose again on the third day, so that he might declare himself to be the life and resurrection of the faithful. Or, indeed, because after his humanity he is seen to be slain, but after his deity, he is immortal and subject to no reproach. Therefore, in the old law, one goat in Leviticus 16 is slain, but the other is not killed, but is led forth into the desert by the work of a man appointed for this purpose. Nevertheless, some interpreters explain it in this way: he is said to be as if he were slain, because, following Saint Chrysostom and Saint Augustine, he has reserved the scars of his death’s wounds as a token of his victory, \&c.

\index{Daniel (Book of)}Furthermore, this Lamb, Christ Jesus our Lord, has seven horns, not that in deed he carries so many horns like a goat of Judea. An horn, as appears by Daniel and by the song of Zacharias in the first chapter of Luke, signifies power and kingdom. The number seven is the number of fullness. It is therefore signified that Christ is endowed with every kind of power, divine, human, imperial, pontifical, royal, briefly, most absolute. In the thirteenth chapter, we shall hear that the beast has taken to him two horns, as it were of the lamb, whereof I shall speak in its place. Daniel in the seventh chapter says, “And rule was given him, and honor and kingdom, that all nations and tongues might worship him, whose rule is an everlasting rule, which shall not perish nor decay at any time.” Now he has seven eyes also. These he expounds, and says, which are the seven spirits of God, sent into the whole world. I showed you before that the seven spirits are called a sevenfold spirit. Here therefore is signified the fullness of the Spirit, which the Lord pours out upon all flesh. Here is signified the universal knowledge of the Son, in whose sight are present whatever things are done in heaven and on earth, openly and privately. For the Spirit of Christ, that unmeasurable force, incomprehensible and most divine, searches and pierces all things; nothing is hidden from his eyes, which view the whole world.

\index{Arethas of Caesarea}And such is Christ, as we have heard described hitherto, whom the patriarchs have before said should come, a victor and triumphant conqueror alone, truly invincible, very man of our own substance, and also our very brother, yet very God nevertheless, of the same substance with the Father and the Holy Ghost, the reconciler, redeemer, and the only salvation of the world: he has suffered for us, and the same has risen again from the dead, and ascended into heaven, having all power in heaven and on earth, who sees all things, communicates his spirit unto men, and is the most faithful keeper and defender of all mankind. This Christ Jesus our Lord came and received; he did not convey or steal it away, but took that book of divine providence, of the judgments of God, of the universal government of all things, that he might open and loose the Seals thereof: that is to say, that he might reveal to us who are redeemed with his blood the judgments of God, and might dispose and order all things in heaven and on earth. Therefore, if we know that the governor of all things is given to us as redeemer, King, Bishop, and our only salvation, who will from henceforth willingly submit himself to his government? And saying we now understand certainly how that under the seal of faith and truth all things are done by Christ, who dares hereafter more curiously inquire of his works and judgments, unto whose credit and government we should now commit all things, in case they were in our power? Nevertheless, we shall observe that the Son does not so receive these things of the Father, that the Father is deprived thereof. For in the fifth chapter of John’s Gospel, the Lord says: “My Father works until now, and I work,” and so forth. Certainly, the Son is called the Word, mouth, and arm of the Father, and so forth, lest after his humanity the Son should seem less than the Father. For very godly Arethas says, where the Lamb received the book from the right hand of him that sits on the Throne, it must be understood on behalf of his humanity, as also that he was slain. For concerning his deity, none of all those things that may worthily be spoken or thought of God is severally assigned to three persons, saving the manner of bringing forth, of him that begets and of him that is begotten, and of him that proceeds, and so forth.

This description of Christ is singular, most excellent, very evangelical, and full of consolation; therefore, it should chiefly be stored in the depths of our hearts. We also see that those who were not afraid to say that, besides the apostolic manner, few things were taught in this book concerning Christ and our redemption, were deceived in their judgment. Let us pray to the Lord that he would deign to illuminate our minds. Amen.
